\chapter{Conclusions and Recommendations}
\section{Summary of Research Findings}
The study documented and experimentally evaluated a text-only complaint-classification component using a controlled synthetic dataset. The 24,000-record corpus was deliberately balanced across three labels and four assigned language groups, with 118 researcher-defined generation categories used for systematic scenario coverage. The model design uses TF--IDF with Linear SVM and a Multinomial Naive Bayes baseline, and separates scenario metadata from the predictive narrative. The original recorded findings, corroborated by reproducibility verification, are: SVM validation/test Macro-F1 remained 1.000000 and baseline Macro-F1 was 0.992077/0.991624. A separate combined-feature candidate tied SVM on development validation; no classification improvement was demonstrated. The group-disjoint cross-validation and fixed validation results are high, although the exact ceiling effects, class-dependent text lengths, and features shared by generation patterns constrain the inference that can be drawn about unfamiliar citizen language. The work offers a reproducible research benchmark and a proposed human-governed complaint processing architecture; it does not establish municipality-wide predictive reliability or measured benefits of live feedback-driven retraining.

\section{Recommendations and Future Work}
\subsection{Data Coverage and Representation Bias}
The 118-category inventory and the synthetic text-generation rules should be audited with fluent Cebuano, Filipino and code-switched language reviewers. The current 8,000/8,000/8,000 class allocation is a controlled experimental design and should not be treated as actual submission prevalence. Future evaluations should include deliberately different reporting-channel styles, rare but meaningful complaints, safety-sensitive reports, varied lengths, misspellings and independent root scenarios. If authorized administrative reports are obtained later, they must be separated from synthetic scenarios by provenance and privacy controls; their collection is not required to state the current experimental findings accurately.

\subsection{Independent Annotation and Feedback Bias}
The presently generated labels have not received measured independent double annotation. A blinded sample spanning the full range of classifications, languages and difficult cases should be reviewed by at least two annotators, with Cohen's kappa and adjudication records reported. Future Human-in-the-Loop correction logs may reflect the judgment of a narrow set of administrators and reporting channels. They should be audited for class prevalence, language, case severity and review consistency before their use in retraining. Model improvement must be tested on independent frozen evaluation examples rather than assumed from the existence of correction records.

\subsection{Generalization, Model Selection and Operational Safety}
Additional testing should use held-out generator root families, unseen complaint intents, minimal pairs changing only the critical context, and realistic class-base-rate scenarios. Regularization, cautious preprocessing, group-aware cross-validation and model versioning should be retained; class weighting or training-only oversampling may be tested \emph{if} feedback data later show meaningful imbalance. Neither SMOTE nor neural dropout is part of the current fitted TF--IDF + linear SVM pipeline. To establish operational readiness, the application also requires authorization, privacy, emergency referral, database and state-transition integrity, administrator audit trails, latency, monitoring and rollback testing. Safety-critical reports must remain accessible to human review even when a text classifier makes an error.

\subsection{Accountable Case Management and Scope Boundaries}
The proposed accepted-case workflow separates Unresolved, Ongoing and Resolved administrative statuses from Valid, Suspicious and Spam model predictions. It recommends a case tracking reference, responsible officer assignment, activity log, policy-specific response target, supervisor notification of overdue events, referral tracking, evidence review before closure and resident reconsideration. These are auditable system-design requirements rather than measured improvements in governmental response or verified legal compliance. Response duties under Republic Act No.~6713 and service-specific processing standards under Republic Act No.~11032 must be interpreted through the applicable Citizen's Charter; no universal deadline for complete incident resolution is assumed \cite{ra6713,ra11032,dcr}.

Optional image attachments, image classification, multimodal fusion, GPS maps and automated authenticity detection are deliberately deferred. A supporting location description can be requested where necessary without requiring GPS. The text-only SVM is not capable of independently deciding whether a physical incident occurred, whether an image is authentic, or whether a complaint was resolved.

\subsection{Evaluation of Proposed Extensions}
The combined word/character feature candidate and grouped sigmoid calibration have now been evaluated on existing synthetic development data only, tying the reference classifier without demonstrated classification improvement. Before field deployment, independently validate transfer and probability reliability and implement/test the class-based review queue; conduct a double-annotation pilot and a representative audit of predicted Valid reports; evaluate a complete case-management prototype; and compare cumulative verified-feedback checkpoints of 100, 250 and 500 only if authentic, quality-checked labels are actually collected. These are experimental checkpoints, not guaranteed improvement thresholds \cite{figueroa2012,ghose2024,fleisig2024}. No prospective feedback-checkpoint or operational HITL results are asserted in the present manuscript.

\section{Concluding Statement}
Controlled synthetic examples make it possible to vary complaint-writing attributes and expose classification weaknesses systematically, while documented EDA and evaluation protect against ungrounded claims. The contribution of the present work is the reproducible investigation of a multilingual three-label complaint-classification task. Further independent human review and operational testing are required before stronger conclusions about deployment or feedback-driven learning can be justified.
